\subsection{An equivalent septic gate in degree fifty-six}

Let $E:y^2+xy=x^3-x^2-2x-1$, let
$\alpha=\sqrt[7]{3}$ be real, and put
\[
 K=\mathbb Q(\zeta_7,\alpha,E[3]),\qquad
 L=\mathbb Q(\alpha,E[3]).
\]
Let $c$ be complex conjugation in this fixed embedding and
write $K'=K^{\langle c\rangle}$ and
$M=L^{\langle c\rangle}$.
We shall verify $[L:\mathbb Q]=112$ and
$D_L=3^{110}7^{124}$ as part of the argument.
The preceding field calculations give $[K:\mathbb Q]=336$ and
$\operatorname{Cl}(K_0)[3]=0$ for the degree-$48$
field $K_0=K^{C_7}=\mathbb Q(\zeta_7,E[3])$.

\begin{proposition}\label{prop:lm-septic-degree56}
Let $z$ be either real root of
\[
 R(Z)=27Z^8-1323Z^6+6174Z^4-117649.
\]
Then
\[
 M=\mathbb Q(z,\alpha),\qquad [M:\mathbb Q]=56,
 \qquad \operatorname{sign}(M)=(2,27),
 \qquad D_M=-3^{55}7^{62}.
\]
Moreover
\[
 \operatorname{Cl}(K)[3]=0
 \quad\Longleftrightarrow\quad
 \operatorname{Cl}(M)[3]=0.
\]
Consequently this explicit degree-$56$ class-group condition
is an equivalent finite gate for the preceding septic
arithmetic-source argument.  Its vanishing is proved below.
\end{proposition}
\begin{proof}
Put $T=\Q(E[3])$ and $F_0=\Q(\sqrt{-7})$.
The CM action embeds $\mathcal O_{F_0}/3=\mathbb F_9$ in
$\operatorname{End}(E[3])$, making $E[3]$ a rank-one module.
The nontrivial automorphism of $F_0$ acts nontrivially on this
endomorphism field, so $F_0\subset T$. Thus
$\operatorname{Gal}(T/F_0)$ is a subgroup of $\mathbb F_9^\times$.
At the good prime three the CM type has one embedding; its
reflex norm on local units is the identity. The local CM reciprocity
formula of \citet[Theorem~5.1]{Lombardo2017} therefore gives the full
scalar group $\mathbb F_9^\times$ as inertia. Consequently
$[T:\Q]=16$, there is one prime over three with $(e,f)=(8,2)$,
and the scalar subgroup acts transitively on the eight nonzero
three-torsion points.

The valuation of $3$ in $T$ is eight, prime to seven, so
$T(\alpha)/T$ has degree seven. This proves $[L:\Q]=112$
and $[K:L]=3$. The latter extension is cyclic and unramified at
three, since the inertia group of $K/\Q$ there has order $56$.
At seven its cubic group lies in decomposition (which has index
two) and then in inertia (whose residue quotient has order two).
Thus both seven-primes are tamely ramified of degree three.
Their residue degrees in $L$ are two, so their contributions to
$N_{L/\Q}\mathfrak d_{K/L}$, the relative discriminant norm, multiply to $7^8$.
The discriminant tower formula gives
\[
 3^{330}7^{380}=D_L^3\,7^8,
 \qquad D_L=3^{110}7^{124}.
\]

First we compare $K'$ and $M$.  The cyclic subgroup
$\operatorname{Gal}(K/L)$ has order three; it fixes
$\alpha$ and $E[3]$ and acts on $\zeta_7$ by a power
of $\zeta_7\mapsto\zeta_7^2$.  It commutes with $c$,
as is checked on these three generators.  Thus $K'/M$
is cyclic of degree three.

The two primes above seven in $K$ are interchanged
by $c$: they are already distinguished by the two
embeddings of $\mu_3$, which lies in $\mathbb Q(E[3])$.
Their contractions to $L$ are likewise interchanged.
Consequently $K'$ and $M$ each have exactly one prime
above seven.  The respective local ramification and
residue indices are $(84,2)$ and $(28,2)$, so that
this prime is totally and tamely ramified in $K'/M$.
At three the cyclic cubic extension is unramified:
inertia in $K/\mathbb Q$ has order $56$, prime to
three.  There is no ramification at other finite
primes, and a cyclic cubic extension cannot ramify
at infinity.  Hence precisely one prime ramifies
in $K'/M$.

Chevalley's ambiguous class number formula, as stated
by \citet[Theorem~1, author p.~1]{Lemmermeyer2013}, now reads
\[
 \bigl|\operatorname{Cl}(K')^{\operatorname{Gal}(K'/M)}\bigr|
 =\frac{h_M}{
 [\mathcal O_M^\times:
  \mathcal O_M^\times\cap N_{K'/M}(K')^\times]}.
\]
The denominator is a power of three, since the cube of every base
unit is a norm. If $3\nmid h_M$,
integrality forces it to have three-adic valuation zero,
so the fixed class group has no three-primary part.
A nonzero finite abelian three-group with an action of
$C_3$ has nonzero fixed elements: on its subgroup
killed by three, $(\sigma-1)^3=0$.  Therefore
$\operatorname{Cl}(K')[3]=0$.

Conversely, the norm from $\operatorname{Cl}(K')$ to
$\operatorname{Cl}(M)$ is surjective.  Indeed, the
cyclic cubic extension ramifies at seven, and hence
intersects the Hilbert class field of $M$ trivially.
Base change of that Hilbert class field and Artin
reciprocity identify the resulting surjection with
the ideal-class norm.  This proves
\[
 \operatorname{Cl}(K')[3]=0
 \quad\Longleftrightarrow\quad
 \operatorname{Cl}(M)[3]=0.
\]

We next compare $K$ and $K'$.  Since two is invertible
on finite three-groups, extension and norm identify
the three-primary class group of $K'$ with the
$c$-invariants of that of $K$.
If $\operatorname{Cl}(K')[3]=0$, then $c$ acts as
$-1$ on $V=\operatorname{Cl}(K)[3]$.
For $g$ in the normal subgroup $C_7$ of
$\operatorname{Gal}(K/\mathbb Q)$, the explicit Kummer
description gives $cgc^{-1}=g^{-1}$.  Since $c$
acts as a scalar on $V$, the actions of $g$ and
$g^{-1}$ on $V$ are equal.  Their common action has
order dividing both two and seven and is therefore
trivial.  Thus $V=V^{C_7}$.

Prime-to-three extension and norm identify these
$C_7$-invariants with $\operatorname{Cl}(K_0)[3]$,
which has already been proved zero.  Hence $V=0$.
The converse follows immediately from the quadratic
norm and extension maps.  Combining this with the preceding equivalence
proves the claimed equivalence.

For the field description, the degree-$16$ Galois
group of $\mathbb Q(E[3])$ acts transitively on its
eight nonzero three-torsion points.  A nonzero real
point is fixed by $c$, so its degree-eight point
field is exactly $\mathbb Q(E[3])^{\langle c\rangle}$.
The previous division-polynomial and subresultant
calculation shows that $z=2y+x$ generates this point
field and has polynomial $R$.
Adjoining the real root $\alpha$ gives degree
$8\cdot7=56$, because the two degrees are coprime.
The resulting field is contained in $M$ and has
its full degree, so they coincide.
There are exactly two nonzero real three-torsion
points and the ordinate coordinate is injective
on these eight points.  Thus the point field has
two real embeddings.  The pure degree-seven field
has one real embedding, giving signature $(2,27)$
for their compositum.

Finally $L/M$ is unramified at finite places.
At seven its two primes are interchanged by $c$.
At three $c$ acts nontrivially on the unramified
quadratic subfield $\mathbb Q_3(\sqrt{-7})$, and
so is outside inertia.  All other finite primes
are already unramified.  The discriminant tower
formula therefore gives $D_L=D_M^2$.
The signature supplies the negative sign, proving
$D_M=-3^{55}7^{62}$.
\end{proof}
